\documentclass[12pt]{article}
\usepackage[utf8]{inputenc}
\usepackage[T1]{fontenc}
\usepackage{hyperref}
\title{San Pedro Quiatoni Zapotec Language Resources for K-12 Curriculum Development}
\author{Compiled by Claude Sonnet 5.5 (claude-sonnet-5-5) for Andr\'as Kornai}
\date{}
\begin{document}
\maketitle

\section{Introduction}
San Pedro Quiatoni Zapotec (ISO 639-3 zpf, Glottocode sanp1261), also called Quiatoni Zapotec or Eastern Tlacolula Zapotec, is spoken in Oaxaca, Mexico. Variants listed include San Pedro Quiatoni, San Pablo Lachiriega, Santo Domingo Narro and Soledad Salinas. The searches found no dedicated dictionary, grammar or school materials.

\section{Primary Educational Resources}
None were found.

\section{Linguistic and Reference Materials}
\begin{itemize}
\item Wikipedia article: \url{https://en.wikipedia.org/wiki/Quiatoni_Zapotec}
\item Endangered Languages Project entry: \url{https://endangeredlanguages.com/lang/10960}
\item AILLA, Zapotec Survey - San Pedro Quiatoni (2008; elicitation recordings and transcriptions): \url{https://islandora-ailla.lib.utexas.edu/islandora/object/ailla%3A249944}
\item AILLA, Dictionary of the Zapotec Language (multi-volume, \'Oscar M\'endez Espinoza). It draws on several Zapotec languages, and it is not confirmed that it covers Quiatoni: \url{https://islandora-ailla.lib.utexas.edu/islandora/object/ailla%3A279600}
\end{itemize}

\section{Religious and Community Resources}
Global Recordings Network lists scripture-related resources (see below). No community organisation was found.

\section{Digital and Audio Resources}
\begin{itemize}
\item Global Recordings Network: \url{https://www.globalrecordings.net/en/language/zpf}
\end{itemize}

\section{Curriculum Development Context}
With about 15,000 speakers the language appears viable, but published educational material is absent from the results. A curriculum would need new materials, building on the AILLA recordings.
\end{document}